\section{Module \texttt{button.c}}

Le module \texttt{button.c} regroupe la création des boutons standards et des boutons radio. Il ajoute au widget GTK de base une logique de thème, d'icône et de raccordement de callbacks.

\subsection{Fonction \texttt{create\_button}}

La fonction \texttt{create\_button} traite plusieurs cas d'usage avec une seule interface : bouton texte, bouton avec mnémonique, bouton avec icône, ou combinaison texte plus icône. La variante de thème est appliquée séparément via une fonction interne dédiée.

\begin{lstlisting}[language=C]
GtkWidget *create_button(const button_config *config) {
    GtkWidget *button;
    const char *label;

    if (config == NULL) {
        return gtk_button_new();
    }

    label = config->label ? config->label : "";
    button = config->use_underline ? gtk_button_new_with_mnemonic(label) : gtk_button_new_with_label(label);
    gtk_button_set_has_frame(GTK_BUTTON(button), config->has_frame);
    gtk_button_set_can_shrink(GTK_BUTTON(button), config->can_shrink);

    if (config->icon_name && config->icon_name[0] != '\0') {
        GtkWidget *content = gtk_box_new(GTK_ORIENTATION_HORIZONTAL, 6);
        GtkWidget *icon = gtk_image_new_from_icon_name(config->icon_name);
        GtkWidget *text = config->use_underline ? gtk_label_new_with_mnemonic(label) : gtk_label_new(label);
        if (config->icon_size > 0) {
            gtk_image_set_pixel_size(GTK_IMAGE(icon), config->icon_size);
        }
        gtk_box_append(GTK_BOX(content), icon);
        gtk_box_append(GTK_BOX(content), text);
        gtk_button_set_child(GTK_BUTTON(button), content);
    }

    apply_theme_variant(button, config->theme_variant);
    apply_widget_style(button, &config->style);

    if (config->on_clicked) {
        g_signal_connect(button, "clicked", G_CALLBACK(config->on_clicked), config->user_data);
    }

    return button;
}
\end{lstlisting}

\subsection{Fonction \texttt{create\_radio\_button}}

Le wrapper radio repose sur \texttt{GtkCheckButton}, comme le recommande GTK4. Le groupement exclusif est obtenu en reliant un bouton au précédent du groupe.

\begin{lstlisting}[language=C]
GtkWidget *create_radio_button(const radio_button_config *config) {
    GtkWidget *button = gtk_check_button_new_with_label(config && config->label ? config->label : "");

    if (config == NULL) {
        return button;
    }

    if (config->group_with != NULL) {
        gtk_check_button_set_group(GTK_CHECK_BUTTON(button), GTK_CHECK_BUTTON(config->group_with));
    }
    gtk_check_button_set_active(GTK_CHECK_BUTTON(button), config->is_active);
    apply_widget_style(button, &config->style);

    if (config->on_toggled) {
        g_signal_connect(button, "toggled", G_CALLBACK(config->on_toggled), config->user_data);
    }

    return button;
}
\end{lstlisting}

\subsection{APIs GTK utilisées}

\begin{itemize}
    \item \textbf{\texttt{gtk\_button\_new\_with\_label()}} : Crée un nouveau widget bouton contenant un label texte. Dans le module, elle est utilisée pour instancier le bouton de base lorsque l'option \texttt{use\_underline} est désactivée, offrant une création simplifiée de boutons textuels standards.
    \item \textbf{\texttt{gtk\_button\_new\_with\_mnemonic()}} : Crée un bouton dont le label contient un mnémonique (raccourci clavier) indiqué par un caractère souligné. Le wrapper y a recours lorsque \texttt{config->use\_underline} est vrai, facilitant ainsi l'accessibilité de l'interface au clavier.
    \item \textbf{\texttt{gtk\_button\_set\_child()}} : Définit le widget enfant unique du bouton. Elle permet au wrapper d'injecter une \texttt{GtkBox} personnalisée contenant à la fois une icône et un texte, dépassant ainsi la limitation des boutons à simple label.
    \item \textbf{\texttt{gtk\_image\_new\_from\_icon\_name()}} : Génère un widget d'image à partir du nom d'une icône présente dans le thème système. Elle est utilisée ici pour charger visuellement l'icône demandée dans la configuration avant son insertion dans le bouton.
    \item \textbf{\texttt{gtk\_check\_button\_set\_group()}} : Lie un bouton de type case à cocher à un groupe existant pour obtenir un comportement de sélection exclusive. Dans \texttt{create\_radio\_button}, cette API transforme un \texttt{GtkCheckButton} standard en un bouton radio fonctionnel.
    \item \textbf{\texttt{g\_signal\_connect()}} : Connecte une fonction de rappel à un signal spécifique d'un objet. Le module l'utilise pour brancher les signaux \texttt{clicked} (pour les boutons) et \texttt{toggled} (pour les radios) aux fonctions définies par le développeur applicatif.
\end{itemize}

\newpage
